\section{\bt{Related Work}{相关工作}}
\label{sec:related}
\paragraph{\bt{Data agents and business semantics}{数据智能体与业务语义}}
\bi{Agent-first data systems identify redundant agent requests and propose a persistent agentic memory that tracks changes to data and metadata~\cite{liu2026agents}; tutorials survey data agents and their difficulty with changing data~\cite{luo2026dataagents} and the management of agentic memory~\cite{li2026agenticmemory}. Enterprise schemas remain hard for text-to-SQL systems~\cite{lei2025spider2,floratou2024nl2sql}, explicit business semantics raise accuracy~\cite{li2023bird,sequeda2024kg}, and interactive benchmarks treat missing business knowledge as a failure to recover from~\cite{huo2026birdinteract}. Business-intelligence platforms maintain term definitions~\cite{jiang2025siriusbi} or generate calculation logic from script history~\cite{weng2025datalab}; AgentSM organizes trajectories as structured semantic memory and leaves the consistency of shared memory to future work~\cite{biswal2026agentsm}, which \system\ provides for metric definitions. Robustness to schema variation~\cite{zhang2025evoschema,furst2025datamodel} is complementary.}
{面向智能体的数据系统指出智能体请求高度重复，并提出跟踪数据与元数据变化的持久智能体记忆~\cite{liu2026agents}；教程综述了数据智能体及其难以适应变化数据的问题~\cite{luo2026dataagents}和智能体记忆的管理~\cite{li2026agenticmemory}。企业表结构对 Text-to-SQL 系统仍然困难~\cite{lei2025spider2,floratou2024nl2sql}，显式业务语义能提高正确率~\cite{li2023bird,sequeda2024kg}，交互式基准把缺失的业务知识当作需要弥补的失败~\cite{huo2026birdinteract}。商业智能平台维护术语定义~\cite{jiang2025siriusbi}或从脚本历史生成计算逻辑~\cite{weng2025datalab}；AgentSM 把轨迹组织为结构化语义记忆，并把共享记忆的一致性留作未来工作~\cite{biswal2026agentsm}，\system\ 为指标定义提供了这种一致性。表结构变化下的鲁棒性~\cite{zhang2025evoschema,furst2025datamodel}与本文互补。}

\paragraph{\bt{Agent memory and experience reuse}{智能体记忆与经验复用}}
\bi{Agent memories induce workflows from successful trajectories~\cite{wang2025awm}, distill insights or strategies~\cite{zhao2024expel,ouyang2026reasoningbank}, evolve context playbooks~\cite{zhang2026ace}, link notes~\cite{xu2025amem}, manage long-term facts~\cite{chhikara2025mem0,packer2023memgpt}, cache plan templates~\cite{zhang2025apc}, and reuse past solutions~\cite{guo2024dsagent}; they hold knowledge about the agent's own tasks and evolve by task feedback. Learning to Share learns which intermediate steps parallel agents add to a memory shared within one task~\cite{fioresi2026learning}. Invalidation contracts let an API server attach table version stamps so that agents evict cached fixes after data drift~\cite{wu2026invalidation}; PlanFence checks before an action that a shared plan was derived from current input versions, without a common snapshot~\cite{chen2026fresh}; MemTX commits writes to a shared memory through snapshot-isolated transactions and revokes derived records along their provenance~\cite{li2026memtx}; and \textsc{MemoRepair} withdraws the descendants of a corrected source and republishes only successors that pass validation~\cite{zhao2026memorepair}. These start from version stamps or explicit revocation events, and none checks stored knowledge against the data it describes. What distinguishes \system\ is that its records describe an external database that changes underneath them: it derives executable conditions from each record's structure, checks them against that database, and binds the evidence to the snapshot a query reads; a failed condition leads to a bounded, regression-tested repair or to invalidation, and an improvement is adopted only when maintained conditions imply its equivalence.}
{智能体记忆从成功轨迹中归纳工作流~\cite{wang2025awm}，提炼洞见或策略~\cite{zhao2024expel,ouyang2026reasoningbank}，演化上下文手册~\cite{zhang2026ace}，链接笔记~\cite{xu2025amem}，管理长期事实~\cite{chhikara2025mem0,packer2023memgpt}，缓存计划模板~\cite{zhang2025apc}，复用既往方案~\cite{guo2024dsagent}；它们保存的是关于智能体自身任务的知识，依据任务反馈演化。Learning to Share 学习并行的智能体把哪些中间步骤加入一个任务内共享的记忆~\cite{fioresi2026learning}。失效契约让 API 服务器附上表版本戳，使智能体在数据漂移后驱逐缓存的修正~\cite{wu2026invalidation}；PlanFence 在动作前检查共享计划是否由输入的当前版本推出，但没有共同快照~\cite{chen2026fresh}；MemTX 通过快照隔离的事务把写入提交到共享记忆，并沿溯源撤销派生记录~\cite{li2026memtx}；\textsc{MemoRepair} 撤下被更正来源的派生条目，只重新发布通过验证的后继~\cite{zhao2026memorepair}。它们从版本戳或显式的撤销事件出发，都不对照记录所描述的数据检查保存的知识。\system\ 的区别在于，它的记录描述的是一个会在其下变化的外部数据库：它从每条记录的结构推出可执行条件，用这个数据库本身检验这些条件，并把证据绑定到查询所读的快照；条件失败则导致有界的、经过回归测试的修复，或使定义失效；改进只有在受维护的条件蕴含其等价时才被采纳。}

\paragraph{\bt{Metric layers, constraints, and validation}{指标层、约束与数据验证}}
\bi{Dimensional modeling conforms dimensions and facts~\cite{kimball2013toolkit}, and measures in SQL attach aggregation semantics to tables so that joins do not double count~\cite{hyde2024measures}. Summarizability research identifies disjointness, completeness, and type compatibility as necessary conditions for correct aggregation~\cite{lenz1997summarizability,mazon2009summarizability,horner2004additivity}. Metric layers declare keys, join properties, and measures by hand~\cite{dbt2026metricflow,databricks2026metricviews}. Integrity checking specializes verification to the updates that can violate a constraint~\cite{nicolas1982integrity,ceri1990constraint}, and derived-relation maintenance first identifies irrelevant updates~\cite{blakeley1989irrelevant}. Soft constraints hold on current data without enforcement~\cite{godfrey2001soft,ilyas2004cords}, and profiling discovers and maintains unique column combinations and dependencies~\cite{abedjan2015profiling,abedjan2014swan,schirmer2019dynfd}. Data-validation systems check declared or inferred constraints on growing data~\cite{schelter2018deequ,breck2019validation,song2021autovalidate,redyuk2021dynamic}, and pipeline tooling runs \code{unique}, \code{not\_null}, and \code{relationships} tests~\cite{dbt2026tests}, expectation suites~\cite{gx2026overview}, and data contracts~\cite{bitol2026odcs}. These checks are declared on tables, and a failure stops a build without saying which consumer depended on it. \system\ derives its checks from each learned definition, including its own filters, proves them sufficient for the definition's aggregate, and lets their outcomes govern the definition on the snapshot the agent reads.}
{维度建模要求各业务过程使用一致的维度和事实~\cite{kimball2013toolkit}；SQL 中的度量把聚合语义附着在表上，使连接不再重复计数~\cite{hyde2024measures}。可汇总性研究指出，不相交、完备和类型相容是正确汇总的必要条件~\cite{lenz1997summarizability,mazon2009summarizability,horner2004additivity}。指标层由人工声明键、连接性质和度量~\cite{dbt2026metricflow,databricks2026metricviews}。完整性检查按可能违反约束的更新执行验证~\cite{nicolas1982integrity,ceri1990constraint}，派生关系维护先识别无关更新~\cite{blakeley1989irrelevant}。软约束在当前数据上成立但不被强制~\cite{godfrey2001soft,ilyas2004cords}，数据画像发现并维护唯一列组合与依赖~\cite{abedjan2015profiling,abedjan2014swan,schirmer2019dynfd}。数据验证系统在增长的数据上检查声明或推断的约束~\cite{schelter2018deequ,breck2019validation,song2021autovalidate,redyuk2021dynamic}，数据管道工具执行 \code{unique}、\code{not\_null} 与 \code{relationships} 测试~\cite{dbt2026tests}、期望集~\cite{gx2026overview}和数据契约~\cite{bitol2026odcs}。这些检查声明在表上，失败时中止构建，却不说明哪个使用方依赖它。\system\ 的检查由每个学到的定义推出，包含定义自己的过滤，证明对定义的聚合是充分的，并让检查结果在智能体读取的快照上约束定义。}

\paragraph{\bt{View maintenance, shared work, and caching}{视图维护、共享计算与缓存}}
\bi{Materialized-view maintenance~\cite{gupta1995maintenance,blakeley1986efficiently} can defer work off the update path~\cite{colby1996deferred} or until a view is referenced~\cite{zhou2007lazy}, the strategy of our full-recheck baseline. Shared execution evaluates common work once~\cite{sellis1988mqo,harizopoulos2005qpipe,giannikis2012shareddb}, and recycling invalidates intermediate results affected by updates~\cite{ivanova2010recycling}. Caches use semantic descriptions~\cite{dar1996semantic}, update-driven invalidation~\cite{garrod2008ferdinand}, input-dependent keys~\cite{gunda2010nectar}, and leases that merge concurrent misses~\cite{nishtala2013memcache}; TxCache keeps cached results transactionally consistent~\cite{ports2010txcache}; build systems verify dependency traces~\cite{mokhov2018build,hammer2014adapton}. \system\ keys check results like build dependencies, coalesces concurrent checks like leased misses, and applies TxCache's principle to check results with a reuse rule proved sound; what it adds are the checks, derived from learned definitions, and the rule that their outcomes govern which definitions agents may use and on which snapshot.}
{物化视图维护~\cite{gupta1995maintenance,blakeley1986efficiently}可以把工作移出更新路径~\cite{colby1996deferred}或推迟到视图被引用时~\cite{zhou2007lazy}，后者正是整定义重查基线的策略。共享执行对公共工作只计算一次~\cite{sellis1988mqo,harizopoulos2005qpipe,giannikis2012shareddb}，中间结果回收使受更新影响的结果失效~\cite{ivanova2010recycling}。缓存使用语义描述~\cite{dar1996semantic}、更新驱动失效~\cite{garrod2008ferdinand}、依赖输入的键~\cite{gunda2010nectar}及合并并发未命中的租约~\cite{nishtala2013memcache}；TxCache 使缓存结果保持事务一致~\cite{ports2010txcache}；构建系统核验依赖轨迹~\cite{mokhov2018build,hammer2014adapton}。\system\ 像构建依赖一样索引检查结果，像带租约的未命中一样合并并发检查，并把 TxCache 的原则用于检查结果，其复用规则被证明可靠；它增加的是由学到的定义推出的检查，以及用检查结果决定智能体可以使用哪些定义、在哪个快照上使用的规则。}
